\section{Related Work}
\label{sec:related}

\paragraph{Memory for language-model agents.}
Memory systems keep a store outside the context window: a paged context
with an archival tier~\cite{packer2023memgpt}, an observation stream with
retrieval and reflection~\cite{park2023generative}, or extracted facts
updated as the conversation goes on~\cite{chhikara2025mem0}, organised by
cognitive architectures into several stores~\cite{sumers2024coala}. A
survey of 435 such works finds the field
``concentrates more heavily on accumulating and retrieving state than on
governing, recovering, or relinquishing it''~\cite{ding2026alwayson}. The
benchmarks that test relinquishing test it on \emph{facts}: answering from
the latest value after an update~\cite{wu2025longmemeval,hu2025memoryagentbench},
not reusing an invalidated preference~\cite{uddin2026memora}, applying a
retrieved update~\cite{chao2026stale}, honouring a withdrawn constraint in
a single decision~\cite{nakayashiki2026stale}. All are graded on a reply;
none involves an obligation whose cue arrives later in a world.

\paragraph{Deferred intentions.}
Two benchmarks test whether an LLM carries out an instruction whose trigger
arrives later, and both touch withdrawal. TriggerBench~\cite{zhang2026triggerbench}
plants a latent constraint in a chat and fires its trigger turns later; in
its negative variants an intermediate turn resolves the constraint, and the
best model stays correctly silent only 44\% of the time.
PM-Bench~\cite{liu2026pmbench} simulates a week of deferred tasks in text
and defines the \emph{update violation}, ``the agent acts on a stale version
of a task after a cancel, override, or reschedule,'' but its week contains
two cancellations, and it reports the hit rate on updated tasks (at most
47\%), not how often the stale version is acted on. Concurrently, Zhao and
Wu~\cite{zhao2026pis} type each intention with a status that code updates
on cancellation and raise PM-Bench Set-F1 from 65 to 83\%, again reporting
misses rather than actions on cancelled tasks. In all three the agent produces a
reply, a reminder or a set of task identifiers, and TriggerBench is scored
by an LLM judge. Smart-home benchmarks let an agent schedule future device
actions and grade them from simulator state~\cite{seo2026simuhome}, but the
scheduled action lives in the device layer and is never withdrawn.
Chameleon~\cite{chameleon2026} studies prospective
memory in a visuomotor policy, without language or cancellation. The
psychology on which we draw separates the two failures these benchmarks
mix: forgetting to act (omission)~\cite{einstein1990normal} and acting on a
finished or cancelled intention (commission)~\cite{scullin2012whoops}.

\paragraph{Revocation, retraction and stopping.}
Recent work reaches revocation from three sides. On the memory side,
Shen et al.~\cite{shen2026revoked} load a revoked policy and its replacement
into five memory systems and find that none enforces revocation by default:
in the two systems that expose the revoked record it is ranked first in
every scenario, and agents take the revoked action in 43\% of trials. Yao et
al.~\cite{yao2026forgetting} show that after the source record is redacted a
revoked preference no longer appears in what the agent says and is still
acted on in 80\% of episodes; retrieval-based and external-memory methods
still leak deleted information~\cite{ren2026gatemem}. On the interaction
side, InterruptBench~\cite{zou2026interrupt} interrupts web agents mid-task,
including with retractions, and scores whether the final intent is achieved,
not how often the retracted requirement is acted on. On the control side,
halt signals sent mid-task are honoured by some models and ignored by
others~\cite{munirathinam2026stop}, and some frontier models subvert a
shutdown that would leave a task incomplete~\cite{schlatter2025shutdown}.
These measure ``stop everything, now''; ours measures ``drop this one thing,
later, while you carry on.''

\paragraph{Grading.}
Our insistence on world-state grading follows evidence that what an agent
says diverges from what it does: rewriting only the reasoning of fixed web
trajectories inflates judge false positives by up to
90\%~\cite{khalifa2026gaming}, and agents misreport their own
progress~\cite{wang2026progressbar}. Contrast-set designs that differ in one perturbation are used to test when agents should
\emph{not} act~\cite{liu2026abstain}; ours vary whether and when an
instruction was withdrawn.

\paragraph{Intentions as commitments.}
In the philosophy of action an intention is a commitment that resists
reconsideration but can be dropped~\cite{bratman1987intention}, the formal
account makes the conditions for dropping it part of the
definition~\cite{cohen1990intention}, and BDI architectures implement this
with an explicit intention structure and a reconsideration
step~\cite{rao1995bdi}. Work joining LLMs to BDI has so far concerned the
interface~\cite{gatti2025chatbdi}, and the case for structured
reconsideration in artificial agents has been made
conceptually~\cite{cariani2024intention}; to our knowledge no study has
measured whether an LLM agent drops an intention. The representation we
find effective is a minimal fragment of a BDI intention structure.
